\section*{第 \ref{ch:b3:identical-particles} 章——全同粒子}

\begin{solution}{exo:b3:identical-particles:1}
(a) 交换两次相当于什么都没做：本征值为 $\pm1$。(b)
全同粒子以对称方式进入 $\hat H$，因此 $\hat H\hat
P_{12} = \hat P_{12}\hat H$：对称类别永不改变。(c)
$\ket a\otimes\ket b$ 与 $\ket b\otimes\ket a$ 将是
\emph{不同}的态，描述同一物理情形——
这种不可观测的区别不应由形式体系携带。(d)
$^1$H（2 个费米子）：玻色子；氘（3）：费米子；$^4$He：玻色子；
$^3$He：费米子；H$_2$（4）：玻色子；光子：玻色子。
\end{solution}

\begin{solution}{exo:b3:identical-particles:2}
(a) $(\ket a\otimes\ket b - \ket b\otimes\ket a)/\sqrt2$. (b)
恒等于零。(c) 对六种置换作交错求和——
即轨道与粒子之间的 $3\times3$ 行列式；
交换两行使行列式变号。(d)
两列相等的行列式为零：两个费米子处于同一
轨道，即两列重复。
\end{solution}

\begin{solution}{exo:b3:identical-particles:3}
(a) $3^2 = 9$。(b) $3$ 个双占据 $+$ $3$ 个对 $= 6$。(c) $3$ 个对。
(d) 六个单粒子态（三个轨道 $\times$ 两种自旋），共
$\binom62 = 15$.
\end{solution}

\begin{solution}{exo:b3:identical-particles:4}
(a) C：$1s^22s^22p^2$；N：$2p^3$；O：$2p^4$；Ne：$2p^6$。(b)
氮的三个 $2p$ 电子各占一个轨道，
自旋平行。(c) 自旋平行迫使空间态反对称，
态，其交换空穴避开了库仑斥力——
自旋排列是以静电能的节省换来的。(d) 氦（及其惰性气体同类）
；氟比氖的满壳层少一个电子，几乎能从任何物质中
夺走一个电子。
\end{solution}

\begin{solution}{exo:b3:identical-particles:5}
(a) 玻色子：$2E_1$；无自旋费米子：$E_1 + 4E_1 = 5E_1$；
电子：$2E_1$（自旋相反者共享 $n = 1$）。(b) 玻色子与电子：
$E_1 + 4E_1 = 5E_1$；无自旋费米子：$E_1 + 9E_1 =
10E_1$。(c) 压缩费米子迫使其占据越来越高的能级：
每次压缩的能量代价是一种向外的压强，任何冷却都无法消除——
这便是物质不可压缩性的雏形。(d)
$\sum_{n=1}^N n^2 \approx N^3/3$：费米海的总能量正比于
$N^3$.
\end{solution}

\begin{solution}{exo:b3:identical-particles:6}
(a) 两项在 $\vect r_1 = \vect r_2$ 处相互抵消。(b) 展开后，
交叉项贡献 $\mp2|\bra a\hat{\vect r}\ket b|^2$：
反对称态具有\emph{更大}的
平均平方间距。(c) 相距越远，斥力能越小：$E_- <
E_+$。(d) 氦的三重态（空间反对称）低于同组态的单重态——
正氦位于仲氦之下，两者相差
\qty{0.8}{eV}，即所避开的斥力。
\end{solution}

\begin{solution}{exo:b3:identical-particles:7}
(a) $1s^2$ 是对称的空间态：自旋须组成
反对称的单态。(b) 分裂 $= 2J$：$J \approx
\qty{0.4}{eV}$。真正的磁偶极—偶极能量在
$a_0$: $\mu_0\mu_{\text{B}}^2/4\pi a_0^3 \approx \qty{4e-4}{eV}$
——小约一千倍：“自旋力”实为静电作用，
只是披上了自旋的外衣。(c) 衰变到 $1\,^1S$ 需要自旋翻转
（单态—三重态）\emph{且}需要 $s \to s$ 跃迁（$\Delta l =
0$）：双重禁戒。(d) $8000$ 秒对 \qty{1.6}{ns}：十二个
数量级，由两条选择定则导出——亚稳性是规则的叠加，
而非力被削弱。
\end{solution}

\begin{solution}{exo:b3:identical-particles:8}
(a) 总交换符号 $=$（自旋部分）$\times (-1)^J$ 必须为
$-1$：三重态（自旋对称）取奇数 $J$，单态取偶数 $J$。
(b) 由重数得 $3:1$。(c) 3/4 的正氢比例
弛豫为仲氢（$J = 0$），释放 $2B \approx \qty{15}{meV}$
每个分子——\emph{多于}液体的潜热：
未经催化的储罐会慢慢自行沸干。工厂则在液化时
借助催化剂将正氢 $\to$ 仲氢。(d) 液氢
火箭技术——每个发射场都继承了
一个核自旋统计问题。
\end{solution}

\begin{solution}{exo:b3:identical-particles:9}
(a) 速率按振幅平方变化：进入已占据模式为 $n + 1$，
反向为 $n$。(b) $10^{10} + 1$ 对 $1$：拥挤的
模式胜出一百亿倍。(c) 一个自发光子落入有利模式，
打破了概率平衡；每个受激光子使 $n$ 增大，
使倾斜加剧：全同光子的雪崩——
增益收窄至单一模式。(d) 在 $\hat a^\dagger\ket n =
\sqrt{n+1}\,\ket{n+1}$ 中：玻色模式就是振子，
而增强因子正是其升降系数。
\end{solution}

\begin{solution}{exo:b3:identical-particles:10}
(a) 每种气体膨胀为两倍体积，$\Delta S =
Nk_{\text{B}}\ln2$（每种气体）。(b) 关闭再重新打开阀门，
将在状态毫无变化的情况下凭空制造熵——
永动式的记账骗局。(c) 若分子有标号，
合并气体多出 $\binom{2N}{N}$ 种“哪一侧”的指派，
相当于 $2Nk_{\text{B}}\ln2$；将所有态数除以 $N!$——
全同粒子的正确计数——恰好消去这一项。
(d) 经典物理需要却无法论证的 $N!$，
正是从远处看到的对称化假设。
\end{solution}

\begin{solution}{exo:b3:identical-particles:11}
(a) $E_{\text{F}} = \hbar^2(3\pi^2n_{\text{e}})^{2/3}/2m_{\text{e}}
\approx \qty{0.16}{MeV}$——约为电子静能的三分之一：
恒星位于相对论区的边缘。(b)
$k_{\text{B}}T \approx \qty{0.9}{keV} \ll E_{\text{F}}$：在费米能标度上，
这颗白热恒星处于深度简并，是“冷”的。
(c) 电子以泡利方式堆叠至
$\qty{0.16}{MeV}$；压缩恒星会抬高每一个被占据的
能级，而这一能量梯度即是与温度无关的压强——
承载重力的是泡利不相容，而非热。(d)
相对论性简并使压强规律趋软，直至引力
在有限质量处占优：即钱德拉塞卡极限，留待
\cref{ch:b3:astrophysics}.
\end{solution}

\begin{solution}{exo:b3:identical-particles:12}
(a) $\hat{\vect S}_1\cdot\hat{\vect S}_2 = \tfrac12(\hat S^2 -
\tfrac32\hbar^2)$：单态上为 $-\tfrac34\hbar^2$，
$+\tfrac14\hbar^2$ 对应三重态——所给哈密顿量
复现 $2J$ 的分裂。(b) 近邻自旋取向一致的晶格：
自发磁化强度——铁磁性。(c) $J/k_{\text{B}}
\approx \qty{1200}{K}$：铁的 \qty{1043}{K} 恰处于
此标度。(d) 偶极–偶极磁作用为 $10^{-4}$ eV：它
将在毫开尔文量级有序化。铁磁性是库仑斥力
借交换几何选择自旋排列。
\end{solution}

\begin{solution}{pb:b3:identical-particles:1}
\textbf{1.} 轨道 $(n, l, m)$，$n^2$ 个分属各 $n$ 层；自旋使之加倍；
泡利不相容限定每个完整标记至多一个电子：$2$、$8$、$18$。
\textbf{2.} 内层电子对原子核的屏蔽因
$l$ 而异：穿透性强的 $s$ 轨道沉至穿透性弱的
$p$ 和 $d$ 轨道之下——能量次序由 $(n, l)$ 决定，而非 $n$
单独决定。
\textbf{3.} 次序为 $1s\,2s\,2p\,3s\,3p\,4s\,3d$；Na：[Ne]$3s^1$；
Cl：[Ne]$3s^23p^5$；K：[Ar]$4s^1$；Fe：[Ar]$4s^23d^6$。
\textbf{4.} 各周期在惰性气体构型处闭合：$2$（$1s$），
$8$（$2s2p$），又为 $8$（$3s3p$——因 $3d$ 高于 $4s$
而需等待），随后为 $18$（$4s3d4p$）。
\textbf{5.} 闭壳层：与下一轨道间隔大，无
易失之电子，亦无处接纳。
\textbf{6.} $Z_{\text{eff}}$ 为核电荷减去
内层（及部分同壳层）电子的屏蔽：钠的价电子
感受到 $11 - 10 \approx 1$；氦的两个电子共享
裸 $Z = 2$，各仅屏蔽对方约三分之一。
\textbf{7.} Li $\to$ Be 上升，在 B 处下降；N $\to$ O 下降：
即两处凹陷。
\textbf{8.} 硼开启 $2p$，其能级高于
$2s$，穿透性更弱，故新增电子更易移去。
\textbf{9.} 氧首次须让两个电子\emph{成对}
占据同一 $p$ 轨道：额外斥力，丧失交换稳定化
——第四个 $2p$ 电子的贡献被打折。
\textbf{10.} 每个电子对另一个的屏蔽都不完全：
拟合 $24.6 = 13.6\,Z_{\text{eff}}^2$ 给出 $Z_{\text{eff}} =
1.34$——约三分之二的电荷被屏蔽。
\textbf{11.} 近乎完整的内核屏蔽（$Z_{\text{eff}}
\approx 1.3$）与外层电子的 $n = 2$：二者共同削弱了
束缚能。
\textbf{12.} 新增电子进入\emph{同一}壳层，彼此
屏蔽很弱，而 $Z$ 不断增大：$Z_{\text{eff}}$ 增大，
整个壳层收缩——氖原子小于锂原子。
\textbf{13.} 惰性气体内核之上有一个易失的 $s$ 电子：电离能
$\sim\qty{5}{eV}$，瞬时生成 $+1$ 离子，给电子剧烈
——族的特性源于同一电子构型。
\textbf{14.} $p^6$ 闭壳层中留有一个空穴：电子亲和能
很大，形成 $-1$ 离子，空穴两两配对成 X$_2$ 分子。
\textbf{15.} $2s$ 与 $2p$ 相距很近：激发 $2s^22p^2 \to$
形成四个半充满的杂化轨道，代价很小，换得两个额外的键——
碳的四价性是对近简并的利用。
\textbf{16.} 四个电子时，成键\emph{与}反键组合同时被填满：
增益相互抵消，不形成共价键
——氦仍为独行者。
\textbf{17.} $d$ 电子埋入 $4s$
外壳之内：外层化学性质由 $s$ 壳层决定，
在十种元素间几乎不变——化学平台。
\textbf{18.} 钪，[Ar]$4s^23d^1$：首个未成对的 $d$
电子。
\textbf{19.} 三个取向一致的自旋给出的磁矩
（$\sim3\mu_{\text{B}}$）；此取向以交换为代价——
库仑能的节省，而非磁力。
\textbf{20.} 若电子略可区分，则所有电子都会
挤向 $1s$（对称化态不再精确相消）：壳层
会渗漏，闭壳层变得模糊，周期律随之瓦解。
\textbf{21.} $v/c \approx Z\alpha = 0.58$：相对论性质量
增大使 $6s$ 轨道收缩并趋于稳定；金的
$5d \to 6s$ 吸收移向蓝端（留下反射的
黄光），汞被收紧的 $6s^2$ 电子对成键极弱，以致
该金属呈液态。
\textbf{22.} 二者能级几乎相当：占据依赖的
次序，$4s$ 先电离，以及 Cr 和 Cu 的 $3d^5$/$3d^{10}$ 借位
现象——半充满与全充满亚壳层以低代价换得
交换稳定性。
\textbf{23.} 化学是电子的事，电子只感受
到原子核的\emph{电荷}：质量仅经
微小的振动（同位素）位移而间接进入。
\textbf{24.} 壳层结构（偏向惰性气体的逻辑）与
相对论稳定化及纯粹的库仑应力在
$Z\alpha \to 1$ 处相抗衡：超重元素化学即二者的拉锯。
\textbf{25.} 轨道的汇总：一次加倍，一个负号，
周期长度 $2, 8, 8, 18$，硼与氧处有凹陷，族从
碱金属到惰性气体，电离能阶梯从 \qty{5.1}{eV} 到
\qty{24.6}{eV}——周期律由推导得出，而非死记硬背。
\end{solution}
